\section{Dataset Description}
The dataset used in this project is the Heart Disease dataset from the UCI Machine Learning Repository~\cite{uci}.
It consists of patient records collected from four different medical centers: Cleveland, Hungarian,
Switzerland, and VA Long Beach. Each subset includes similar clinical attributes that describe the health
status of patients and the presence or absence of heart disease.

In this work, we use the integrated version of the dataset, which contains approximately
\textbf{1000} instances and \textbf{13} main attributes, along with a binary target variable indicating
whether the patient is diagnosed with heart disease (\texttt{target = 1}) or not (\texttt{target = 0}).
The attributes can be grouped into three categories:

\begin{itemize}
	\item \textbf{Numeric (continuous) features}, such as:
	\begin{itemize}
		\item \texttt{age}: age in years
		\item \texttt{trestbps}: resting blood pressure (in mm Hg)
		\item \texttt{chol}: serum cholesterol (in mg/dl)
		\item \texttt{thalach}: maximum heart rate achieved
		\item \texttt{oldpeak}: ST depression induced by exercise relative to rest
	\end{itemize}
	\item \textbf{Binary or categorical features encoded as integers}, for example:
	\begin{itemize}
		\item \texttt{sex}: sex (0 = female, 1 = male)
		\item \texttt{cp}: chest pain type (0--3)
		\item \texttt{fbs}: fasting blood sugar $>$ 120 mg/dl (0 = false, 1 = true)
		\item \texttt{restecg}: resting electrocardiographic results
		\item \texttt{exang}: exercise induced angina (0 = no, 1 = yes)
		\item \texttt{slope}: the slope of the peak exercise ST segment
		\item \texttt{ca}: number of major vessels (0--3) colored by fluoroscopy
		\item \texttt{thal}: thalassemia (0–3, encoded)
	\end{itemize}
	\item \textbf{Target variable}:
	\begin{itemize}
		\item \texttt{target}: presence of heart disease (0 = no, 1 = yes)
	\end{itemize}
\end{itemize}

In the implementation, the dataset is loaded as a CSV file using the \texttt{pandas} library.
Basic inspection of the data shape, column types, and missing values is performed to verify
that the dataset is correctly integrated and ready for preprocessing.